% TOP_PROBLEM: {"id": 3900005, "title": "Triangulating a hypercube", "category_id": 6, "category": {"id": 6, "name": "geometry", "display_name": "Geometry", "description": "Euclidean and non-Euclidean geometry, geometric structures.", "slug": "geometry", "order_index": 6, "created_at": "2026-07-31T15:26:25.671Z"}, "status": "partially_solved", "problem_number": "AMR-038-0005", "set_id": 13, "statusQualification": "The vertex-only and Steiner-vertex conventions are stated separately because the short catalogue wording does not distinguish them."}
% FIRST_POSTED: 2026-10-01
% ENTRY_KIND: statement_only
% UnsolvedMath ID 3900005; source catalogue code AMR-038-0005
% !TeX program = lualatex
% Definition and sources only; this file is not an LLM solution attempt.
\documentclass[11pt,a4paper]{article}
\usepackage[margin=25mm]{geometry}
\usepackage{fontspec}
\setmainfont{lmroman10-regular.otf}[BoldFont=lmroman10-bold.otf,
 ItalicFont=lmroman10-italic.otf,BoldItalicFont=lmroman10-bolditalic.otf]
\usepackage{amsmath,amssymb,mathtools}
\usepackage{unicode-math}
\setmathfont{latinmodern-math.otf}
\AtBeginDocument{\let\setminus\smallsetminus}
\usepackage{xurl,hyperref}
\hypersetup{colorlinks=true,urlcolor=blue}
\setcounter{secnumdepth}{0}
\setlength{\parindent}{0pt}
\setlength{\parskip}{5pt}
\setlength{\emergencystretch}{3em}
\begin{document}
\section*{Triangulating a hypercube}\label{3900005}
{\small UnsolvedMath ID 3900005 · AMR-038-0005 · Geometry · AMR Open Problem Lists}
\subsection{Definitions and mathematical statement}
Let $C_d=[0,1]^d$. A triangulation of $C_d$ is a finite collection of
nondegenerate closed $d$-simplices whose union is $C_d$, such that any two
intersect in a common face or not at all. Its size is the number of its
$d$-dimensional simplices, not the total number of faces.

In the standard vertex-triangulation problem every simplex vertex belongs
to $\{0,1\}^d$. Define
\[
 \tau_d=\min\{\#\mathcal T:\mathcal T\text{ triangulates }C_d
                 \text{ using only cube vertices}\}.
\]
Determine $\tau_d$ and its asymptotic growth as $d\to\infty$, including
sharp exponential improvements over the elementary $d!$ subdivision.
That subdivision orders the $d$ coordinates in every possible way and
uses one simplex for each permutation.

If extra vertices are permitted, define $\tau_d^{\mathrm{St}}$ by the same
minimum with arbitrary vertices in $C_d$. These are Steiner vertices, and
$\tau_d^{\mathrm{St}}\le\tau_d$. The two conventions should be distinguished
when reporting a minimum or an asymptotic bound. A dissection allowing
nonmatching faces, or a demand that all simplices be congruent, is also a
different problem. No congruence condition is imposed here.
\subsection{Short English statement}
What is the smallest triangulation of a high-dimensional cube, and how
does its size grow with dimension? Distinguish cube-vertex triangulations
from those allowed to insert extra vertices.
\subsection{Sources}
\begin{itemize}
\item[S1] ULAM.AI and the UnsolvedMath Contributors, supplied problems.json, record 3900005 (AMR-038-0005), AMR Open Problem Lists. Catalogue identity and source transcription; CC BY 4.0.
\url{https://huggingface.co/datasets/ulamai/UnsolvedMath}.
\item[S2] Eppstein - The Geometry Junkyard: Open Problems, Triangulating a hypercube. Original source indexed by AMR-038-0005.
\url{https://ics.uci.edu/~eppstein/junkyard/open.html}.
\item[S3] David Eppstein, Geometry Junkyard, Cube triangulation.
\url{https://ics.uci.edu/~eppstein/junkyard/cube-triangulation.html}.
\end{itemize}
\end{document}
